\chapter{Rachunek różniczkowy na podrozmaitościach}
\label{ma2:ch:manifold-calculus}

W rozdziale \ref{ma2:ch:submanifolds} opisaliśmy podrozmaitości
przestrzeni euklidesowej i ich przestrzenie styczne. Teraz
rozpatrzymy funkcje i odwzorowania określone na tych zbiorach.
Korzystamy z podejścia przez parametryzacje; ogólny rachunek
na rozmaitościach abstrakcyjnych jest przedmiotem geometrii
różniczkowej. W tym rozdziale podrozmaitości nie mają brzegu
i są klasy $C^\infty$, chyba że jawnie przyjmujemy słabszą
regularność. Wszystkie przestrzenie euklidesowe są rzeczywiste.

\section{Odwzorowania i funkcje różniczkowalne}

Niech $M\subseteq\R^m$ i $N\subseteq\R^n$ będą
podrozmaitościami wymiarów odpowiednio $k$ i $l$.

\begin{definicja}
\label{ma2:def:smooth-map-manifolds}
Niech $F:M\to N$ będzie ciągłe. Mówimy, że $F$ jest klasy
$C^r$ w otoczeniu $p\in M$, gdy dla pewnych parametryzacji
lokalnych $\sigma:U\to M$, $\tau:V\to N$,
$\sigma(a)=p$, $\tau(b)=F(p)$, po odpowiednim zmniejszeniu
$U$ mamy $F(\sigma(U))\subseteq\tau(V)$ oraz
\[
 f=\tau^{-1}\circ F\circ\sigma:U\to V
\]
jest klasy $C^r$. Jeżeli warunek zachodzi w otoczeniu każdego
punktu, mówimy, że $F$ jest klasy $C^r$ na $M$.
Dla $r=\infty$ mówimy o odwzorowaniu gładkim.
Funkcja rzeczywista jest szczególnym przypadkiem $N=\R$.
\end{definicja}

Przy podrozmaitościach klasy $C^s$ definicję stosujemy dla
$1\le r\le s$. Wymaganie ciągłości zapewnia, że można
zmniejszyć dziedzinę parametryzacji źródłowej tak, aby
obraz zawierał się w wybranym otoczeniu parametryzacyjnym
punktu $F(p)$.

\begin{stwierdzenie}
\label{ma2:prop:smooth-map-characterizations}
Definicja \ref{ma2:def:smooth-map-manifolds} nie zależy od
wyboru parametryzacji. Ponadto równoważne są warunki:
\begin{enumerate}[label=(\roman*)]
\item $F$ jest klasy $C^r$ w otoczeniu $p$;
\item $F\circ\sigma:U\to\R^n$ jest klasy $C^r$
w otoczeniu $a$, dla parametryzacji $\sigma(a)=p$;
\item istnieje otwarte otoczenie $W\subseteq\R^m$ punktu $p$
i odwzorowanie $\widetilde F:W\to\R^n$ klasy $C^r$,
takie że $\widetilde F|_{W\cap M}=F|_{W\cap M}$.
\end{enumerate}
Nie wymagamy, aby $\widetilde F$ przyjmowało poza $M$
wartości w $N$.
\end{stwierdzenie}
\begin{proof}
Najpierw zauważmy, że odwrotność dowolnej parametryzacji
lokalnej $\sigma:U\to M$ ma lokalne przedłużenie klasy
$C^s$ do otoczenia w $\R^m$. Istotnie, $D\sigma(a)$ ma
rząd $k$, więc można wybrać rzut współrzędnościowy
$P:\R^m\to\R^k$, dla którego $P D\sigma(a)$ jest
izomorfizmem. Twierdzenie \ref{ma2:thm:inverse-function}
zastosowane do $P\circ\sigma$ daje lokalną odwrotność $H$.
Na dostatecznie małym względnie otwartym fragmencie $M$
mamy $\sigma^{-1}=H\circ P$. Prawa strona jest określona
na otwartym zbiorze w $\R^m$. Ten fragment można zapisać
lokalnie jako $M\cap W$, ponieważ obraz parametryzacji
jest względnie otwarty. Dla $k=0$ odwrotność jest lokalnie
odwzorowaniem stałym do $\R^0$.

W konsekwencji każde odwzorowanie przejścia między
parametryzacjami jest klasy $C^s$: lokalnie ma postać
$H\circ P\circ\sigma_1$. Jego odwrotność otrzymujemy
zamieniając parametryzacje. Po zmianie parametryzacji
źródłowej i docelowej reprezentacja $f$ zmienia się na
\[
 (\tau_1^{-1}\circ\tau)\circ f\circ
 (\sigma^{-1}\circ\sigma_1).
\]
Oba odwzorowania przejścia są klasy $C^s$, więc klasa
$C^r$ jest zachowana w obu kierunkach.

Jeżeli zachodzi (i), to $F\circ\sigma=\tau\circ f$,
co daje (ii). Odwrotnie, dla parametryzacji $\tau$ wokół
$F(p)$ stosujemy skonstruowane wyżej przedłużenie jej
odwrotności i składamy je z $F\circ\sigma$. Po zmniejszeniu
otoczenia, co jest możliwe dzięki ciągłości $F$, otrzymujemy
(ii)$\Rightarrow$(i).
Przy (ii) definiujemy lokalnie
$\widetilde F=(F\circ\sigma)\circ H\circ P$,
gdzie $H\circ P$ przedłuża $\sigma^{-1}$.
Zmniejszamy $W$, aby $H(P(W))\subseteq U$ i
$M\cap W$ leżał w rozpatrywanym obrazie $\sigma$.
To daje (iii). Z (iii) wynika (ii) przez złożenie
$\widetilde F\circ\sigma$. Wszystkie złożenia mają
żądaną regularność na mocy reguł rachunku różniczkowego.
\end{proof}

W szczególności ograniczenie funkcji gładkiej z otoczenia
podrozmaitości jest gładkie na podrozmaitości. Suma,
iloczyn i złożenie gładkich funkcji, przy odpowiednich
dziedzinach i przeciwdziedzinach, są gładkie: we
współrzędnych lokalnych są to zwykłe działania na
funkcjach gładkich. Gładkie funkcje rzeczywiste na $M$
tworzą algebrę $C^\infty(M)$ nad $\R$ z działaniami
punktowymi i jedynką będącą funkcją stałą $1$.

\section{Różniczka odwzorowania}

\begin{twierdzenie}
\label{ma2:thm:manifold-differential}
Jeżeli $F:M\to N$ jest klasy $C^1$ w otoczeniu $p$,
to istnieje dokładnie jedno odwzorowanie liniowe
\[
 \mathrm dF_p:T_pM\to T_{F(p)}N
\]
takie, że dla każdej krzywej klasy $C^1$ w $M$
z $\gamma(0)=p$ zachodzi
\[
 \mathrm dF_p(\gamma'(0))=(F\circ\gamma)'(0).
\]
Jeżeli $\widetilde F$ jest lokalnym przedłużeniem klasy
$C^1$ jak w stwierdzeniu
\ref{ma2:prop:smooth-map-characterizations}, to
$\mathrm dF_p=D\widetilde F(p)|_{T_pM}$.
\end{twierdzenie}
\begin{proof}
Według stwierdzenia \ref{ma2:prop:smooth-map-characterizations}
istnieje $\widetilde F$. Dla każdej krzywej $\gamma$
z definicji, po zmniejszeniu przedziału parametrów,
$F\circ\gamma=\widetilde F\circ\gamma$ jest klasy $C^1$
i na mocy reguły łańcuchowej ma w zerze pochodną
$D\widetilde F(p)\gamma'(0)$. Krzywa $F\circ\gamma$
przyjmuje wartości w $N$, więc jej prędkość należy do
$T_{F(p)}N$. Każdy $v\in T_pM$ jest prędkością
takiej krzywej na mocy definicji i twierdzenia
\ref{ma2:thm:tangent-descriptions}. Ograniczenie
$D\widetilde F(p)|_{T_pM}$ przyjmuje zatem wartości w
$T_{F(p)}N$, jest liniowe i spełnia żądaną tożsamość.
Jeżeli dwa odwzorowania spełniają tę tożsamość, ich
wartości na dowolnym $v$ są równe pochodnej
$F\circ\gamma$ dla reprezentującej go krzywej.
Jednoznaczność dowodzi również niezależności od
przedłużenia.
\end{proof}

\begin{wniosek}[Postać współrzędnościowa]
\label{ma2:cor:manifold-coordinate-differential}
Dla parametryzacji $\sigma(a)=p$, $\tau(b)=F(p)$ i
$f=\tau^{-1}\circ F\circ\sigma$ mamy
\[
 \mathrm dF_p=D\tau(b)\circ Df(a)\circ
              (D\sigma(a))^{-1}.
\]
Odwrotność $D\sigma(a)$ jest tu rozumiana jako odwrotność
izomorfizmu $\R^k\to T_pM$, nie jako odwrotność operatora
o przeciwdziedzinie $\R^m$.
\end{wniosek}
\begin{proof}
Dla $h\in\R^k$ stosujemy twierdzenie
\ref{ma2:thm:manifold-differential} do krzywej
$\gamma(t)=\sigma(a+th)$. Ponieważ
$F\circ\gamma(t)=\tau(f(a+th))$, reguła łańcuchowa
daje $\mathrm dF_p D\sigma(a)h=D\tau(b)Df(a)h$.
Pochodna $D\sigma(a)$ jest izomorfizmem na $T_pM$ na
mocy twierdzenia \ref{ma2:thm:tangent-descriptions}.
Stąd wynika wzór.
\end{proof}

\begin{stwierdzenie}[Reguła łańcuchowa]
\label{ma2:prop:manifold-chain}
Jeżeli $F:M\to N$ oraz $G:N\to S$ są klasy $C^1$,
to $G\circ F$ jest klasy $C^1$ i
\[
 \mathrm d(G\circ F)_p=\mathrm dG_{F(p)}\circ\mathrm dF_p.
\]
Ponadto $\mathrm d(\id_M)_p=I_{T_pM}$. Jeżeli $F$ jest
dyfeomorfizmem klasy $C^1$, to $\mathrm dF_p$ jest
izomorfizmem, którego odwrotnością jest
$\mathrm d(F^{-1})_{F(p)}$.
\end{stwierdzenie}
\begin{proof}
Regularność złożenia wynika z jego reprezentacji w
parametryzacjach. Niech $\gamma$ reprezentuje $v\in T_pM$.
Krzywa $F\circ\gamma$ reprezentuje $\mathrm dF_p v$,
więc dwukrotne użycie twierdzenia
\ref{ma2:thm:manifold-differential} daje
\[
 ((G\circ F)\circ\gamma)'(0)
 =\mathrm dG_{F(p)}((F\circ\gamma)'(0))
 =\mathrm dG_{F(p)}(\mathrm dF_p v).
\]
Dla identyczności pochodna krzywej się nie zmienia.
Stosując wzór do obu złożeń $F^{-1}\circ F$ oraz
$F\circ F^{-1}$, otrzymujemy obie tożsamości odwrotności.
\end{proof}

\paragraph{Lokalne postacie.}
Różniczkę nazywamy różnowartościową (odpowiednio
surjektywną), gdy taki jest jej operator liniowy.
Z wniosku \ref{ma2:cor:manifold-coordinate-differential}
wynika, że jej rząd jest równy rzędowi $Df(a)$.
W szczególności lokalne postacie immersji, submersji
i odwzorowań stałego rzędu z rozdziału 20 stosują się
do reprezentacji współrzędnościowej $f$.
Dla stałego rzędu $j$ w otoczeniu punktu po odpowiednich
zmianach współrzędnych postać ma postać
$(u_1,\ldots,u_k)\mapsto(u_1,\ldots,u_j,0,\ldots,0)$,
przy współrzędnych ze środkiem w rozpatrywanych punktach.
Jeśli $k=l$ i $\mathrm dF_p$ jest izomorfizmem, twierdzenie
\ref{ma2:thm:inverse-function} zastosowane do $f$ daje
lokalny dyfeomorfizm $F$. Jego lokalna odwrotność ma
postać $\sigma\circ f^{-1}\circ\tau^{-1}$.
Nie zakładamy w tych wnioskach globalnej odwracalności.

\section{Funkcje rzeczywiste i gradient styczny}

Dla funkcji $f:M\to\R$ utożsamiamy $T_{f(p)}\R$ z $\R$.
Różniczka $\mathrm df_p$ jest wtedy elementem przestrzeni
$\mathcal L(T_pM,\R)$, nazywanym kowektorem. Od niego
odróżniamy wektor gradientu, zależny od iloczynu skalarnego.

\begin{definicja}
\label{ma2:def:tangential-gradient}
Gradientem stycznym funkcji klasy $C^1$ w $p$ nazywamy
jedyny wektor $\nabla_M f(p)\in T_pM$ taki, że
\[
 \mathrm df_p(v)=\langle\nabla_M f(p),v\rangle
 \qquad(v\in T_pM).
\]
Iloczyn skalarny pochodzi z otaczającej przestrzeni
euklidesowej.
\end{definicja}

\begin{stwierdzenie}
\label{ma2:prop:gradient-projection}
Dla lokalnego przedłużenia $\widetilde f$ klasy $C^1$
mamy
\[
 \nabla_M f(p)=P_{T_pM}\nabla\widetilde f(p),
\]
gdzie $P_{T_pM}$ oznacza rzut ortogonalny.
Jeśli $\sigma(a)=p$, $A=D\sigma(a)$ oraz
$B=A^{\mathsf T}A$, to
\[
 \nabla_M f(p)=A B^{-1}\nabla(f\circ\sigma)(a)
 \qquad(k>0).
\]
Macierz $B$ jest macierzą Grama wektorów bazowych
$\partial_1\sigma(a),\ldots,\partial_k\sigma(a)$.
\end{stwierdzenie}
\begin{proof}
Dla $v\in T_pM$ z twierdzenia
\ref{ma2:thm:manifold-differential} wynika
\[
 \mathrm df_p(v)=D\widetilde f(p)v
 =\langle\nabla\widetilde f(p),v\rangle
 =\langle P_{T_pM}\nabla\widetilde f(p),v\rangle.
\]
Ostatnia równość wynika z ortogonalności składowej
normalnej do $v$. To dowodzi istnienia gradientu;
jeśli dwa wektory styczne dają ten sam funkcjonał,
ich różnica $w$ spełnia $\langle w,w\rangle=0$,
a więc jest zerowa.

Ponieważ $A$ jest różnowartościowe, dla $c\ne0$
mamy $c^{\mathsf T}Bc=\|Ac\|^2>0$. Macierz $B$ jest
zatem odwracalna. Zapisujemy gradient styczny
jednoznacznie jako $Ac$ i dla każdego $h\in\R^k$
otrzymujemy
\[
 D(f\circ\sigma)(a)h
 =\mathrm df_p(Ah)=\langle Ac,Ah\rangle
 =\langle Bc,h\rangle.
\]
Stąd $Bc=\nabla(f\circ\sigma)(a)$ i żądany wzór.
W wymiarze zerowym przestrzeń styczna i gradient są
zerowe.
\end{proof}

\begin{przyklad}[Funkcja wysokości na sferze]
\label{ma2:ex:sphere-height-gradient}
Na sferze $S^2=\{p\in\R^3:\|p\|=1\}$ rozpatrujemy
$f(p)=p_3$. Przestrzeń styczna to $p^\perp$,
na mocy twierdzenia \ref{ma2:thm:tangent-descriptions}
zastosowanego do równania $\|p\|^2-1=0$.
Gradient przedłużenia $\widetilde f(x)=x_3$ to $e_3$.
Zatem
\[
 \nabla_{S^2} f(p)=e_3-\langle e_3,p\rangle p=e_3-p_3p.
\]
Jego kwadrat normy jest równy $1-p_3^2$, ponieważ
$\|p\|=1$. Gradient znika dokładnie na obu biegunach.
W pozostałych punktach wskazuje kierunek najszybszego
wzrostu wysokości wśród jednostkowych wektorów stycznych:
wynika to z nierówności Cauchy'ego--Schwarza i podstawienia
$v=\nabla_{S^2}f(p)/\|\nabla_{S^2}f(p)\|$.
\end{przyklad}

\section{Punkty krytyczne i ekstrema}

\begin{twierdzenie}
\label{ma2:thm:manifold-critical}
Jeżeli $f:M\to\R$ jest klasy $C^1$ i $p$ jest jej
ekstremum lokalnym na $M$, to $\mathrm df_p=0$.
Równoważnie $\nabla_M f(p)=0$. Punkt, w którym zachodzi
ten warunek, nazywamy punktem krytycznym.
\end{twierdzenie}
\begin{proof}
Niech $\sigma(a)=p$ będzie parametryzacją lokalną.
Homeomorfizm $\sigma$ na względnie otwarty fragment $M$
przenosi lokalne ekstremum w $p$ na lokalne ekstremum
$f\circ\sigma$ w otwartym zbiorze $U\subseteq\R^k$.
Twierdzenie \ref{ma2:thm:extremum-necessary} daje
$D(f\circ\sigma)(a)=0$. Z reguły łańcuchowej
$D(f\circ\sigma)(a)=\mathrm df_p\circ D\sigma(a)$;
surjektywność $D\sigma(a):\R^k\to T_pM$ daje
$\mathrm df_p=0$. Równoważność z zanikiem gradientu
wynika z definicji \ref{ma2:def:tangential-gradient}.
\end{proof}

\begin{twierdzenie}[Druga różniczka w punkcie krytycznym]
\label{ma2:thm:critical-hessian-manifold}
Niech $f$ będzie klasy $C^2$ na podrozmaitości klasy
co najmniej $C^2$, a $p$ niech będzie punktem krytycznym.
Wtedy wzór
\[
 H_p(v,w)=D^2(f\circ\sigma)(a)
       [(D\sigma(a))^{-1}v,(D\sigma(a))^{-1}w]
\]
określa symetryczną formę dwuliniową na $T_pM$,
niezależną od parametryzacji $\sigma(a)=p$.
Jeżeli $p$ jest minimum lokalnym, to $H_p(v,v)\ge0$.
Jeżeli $H_p(v,v)>0$ dla wszystkich $v\ne0$, to $p$ jest
ścisłym minimum lokalnym. Dla maksimum znaki są przeciwne.
\end{twierdzenie}
\begin{proof}
Niech $\sigma_1(a_1)=p$ i
$\theta=\sigma^{-1}\circ\sigma_1$ lokalnie wokół $a_1$.
Z reguły drugiej pochodnej złożenia
(stwierdzenie \ref{ma2:prop:second-chain}) otrzymujemy
\[
 \begin{aligned}
 D^2(f\circ\sigma_1)(a_1)[h,k]
 &=
 D^2(f\circ\sigma)(a)[D\theta(a_1)h,D\theta(a_1)k]\\
 &\quad+D(f\circ\sigma)(a)D^2\theta(a_1)[h,k].
 \end{aligned}
\]
Ostatni składnik znika, bo
$D(f\circ\sigma)(a)=\mathrm df_p D\sigma(a)=0$.
Z $\sigma_1=\sigma\circ\theta$ dostajemy
$D\sigma_1(a_1)=D\sigma(a)D\theta(a_1)$, co po
podstawieniu dowodzi niezależności formy. Symetria
wynika z twierdzenia Schwarza dla funkcji klasy $C^2$.

Warunek konieczny otrzymujemy z
twierdzenia \ref{ma2:thm:extremum-necessary}
dla $f\circ\sigma$. Przy dodatniej określoności
$H_p$ forma $D^2(f\circ\sigma)(a)$ jest dodatnio określona
w $\R^k$, ponieważ $D\sigma(a)$ jest różnowartościowe.
Jeżeli $k>0$, zwartość sfery i ciągłość formy dają
jednostajną dodatnią dolną granicę na sferze; stosujemy
twierdzenie \ref{ma2:thm:extremum-sufficient}.
Ścisłe minimum przenosi się przez homeomorfizm
$\sigma$ na $M$. Przy $k=0$ punkt jest izolowany i
ścisłe minimum oraz maksimum lokalne zachodzą
w sensie definicji. Dla maksimum rozpatrujemy $-f$.
\end{proof}

Poza punktem krytycznym zwykłe drugie pochodne w mapach
nie transformują się na ogół jak forma dwuliniowa:
nie musi znikać składnik zawierający $D^2\theta$.
Na przykład na $M=\R$ dla $f(x)=x$ i parametryzacji
$\sigma(t)=t+t^2$ ograniczonej do $(-1/4,1/4)$
mamy $(f\circ\sigma)''(0)=2$, podczas gdy dla
parametryzacji identycznościowej druga pochodna wynosi
zero. Obie parametryzacje mają tę samą pochodną w zerze.
Nie należy więc bez dodatkowej struktury nazywać
macierzy drugich pochodnych w dowolnej mapie
niezależnym od mapy Hesjanem.

\begin{wniosek}[Związek z mnożnikami Lagrange'a]
\label{ma2:cor:manifold-hessian-lagrange}
Niech lokalnie $M=\{x:G(x)=0\}$, gdzie $G$ jest klasy
$C^2$ i $DG(p):\R^m\to\R^{m-k}$ jest surjektywne.
Dla lokalnego przedłużenia $\widetilde f$ klasy $C^2$
w punkcie krytycznym $p$ istnieje jedyny funkcjonał
$\lambda:\R^{m-k}\to\R$ taki, że
$D\widetilde f(p)=\lambda DG(p)$. Wtedy
\[
 H_p(v,w)=D^2\widetilde f(p)[v,w]
                 -\lambda(D^2G(p)[v,w]).
\]
\end{wniosek}
\begin{proof}
Ponieważ $T_pM=\ker DG(p)$, twierdzenie
\ref{ma2:thm:manifold-differential} i krytyczność dają
$D\widetilde f(p)|_{\ker DG(p)}=0$.
Faktoryzacja przez surjekcję i jednoznaczność wynikają
z argumentu liniowego w dowodzie twierdzenia
\ref{ma2:thm:lagrange-multipliers} (nie trzeba zakładać
ekstremum). Dla parametryzacji $\sigma(a)=p$ funkcja
$L=\widetilde f-\lambda\circ G$ spełnia
$L\circ\sigma=f\circ\sigma$ i $DL(p)=0$.
Reguła drugiej pochodnej złożenia daje
$D^2(f\circ\sigma)(a)[h,k]=D^2L(p)[D\sigma(a)h,D\sigma(a)k]$.
Po podstawieniu do definicji $H_p$ otrzymujemy tezę.
\end{proof}

W przykładzie \ref{ma2:ex:sphere-height-gradient} bieguny
$p_\pm=\pm e_3$ są punktami krytycznymi.
Dla $G(x)=\|x\|^2-1$ mnożnik spełnia
$e_3=2\lambda p_\pm$, więc $\lambda=\pm1/2$.
Ponieważ $D^2\widetilde f=0$, a $D^2G=2I$ w zapisie
macierzowym, na $T_{p_\pm}S^2$ forma $H_{p_\pm}(v,w)$
jest równa $\mp\langle v,w\rangle$. Otrzymujemy ścisłe
maksimum na biegunie północnym i ścisłe minimum na
południowym. Są to jedyne ekstrema globalne, gdyż
$-1\le p_3\le1$ z równością wyłącznie w odpowiednim
biegunie.

% Drugi blok rozdziału: orientacja, macierz Grama i element
% pola, ilustracje powierzchni oraz powiązanie z całkowaniem.
